% wherelist:
%   A compact IEEE-like "where" list for explaining symbols after an equation.
%   The word "where" is printed automatically.
%
% Usage:
%
%   \begin{equation}
%   ...
%   \end{equation}
%   \begin{wherelist}
%     \item[$p_i$] is the error rate of the i-th error position;
%     \item[$e_i$] is 1 if the i-th error occurred and 0 otherwise.
%   \end{wherelist}
%
% Optional arguments:
%
%   \begin{wherelist}[<left indent>][<gap between label and explanation>]
%
% Defaults:
%
%   \begin{wherelist}              % equivalent to [\parindent][1em]
%   \begin{wherelist}[1em]         % equivalent to [1em][1em]
%   \begin{wherelist}[1em][0.75em] % custom indent and custom label gap
%
\NewDocumentEnvironment{wherelist}{O{\parindent} O{1em} +b}{%
  \begingroup

  % Local spacing settings.
  %
  % This tightens the space between the preceding displayed equation and
  % the printed word "where". Make this less negative if the gap is too tight,
  % or more negative if the gap is still too large.
  \def\wherelistEquationToWhereCorrection{-0.45\baselineskip}%

  % This removes the extra paragraph/list gap between "where" and the items.
  \def\wherelistWhereToItemsCorrection{-\parskip}%

  % Local readable aliases for TeX scratch dimensions.
  \dimendef\wherelistLabelWidth=0
  \dimendef\wherelistCandidateWidth=2

  \wherelistLabelWidth=0pt
  \wherelistCandidateWidth=0pt

  % First pass: measure the widest optional label in \item[...].
  {%
    \RenewDocumentCommand{\item}{o}{%
      \IfNoValueTF{##1}{%
        \wherelistCandidateWidth=0pt
      }{%
        \settowidth{\wherelistCandidateWidth}{##1}%
      }%
      \ifdim\wherelistCandidateWidth>\wherelistLabelWidth
        \global\wherelistLabelWidth=\wherelistCandidateWidth
      \fi
      \ignorespaces
    }%
    \setbox0=\vbox{%
      \hsize=\maxdimen
      \parindent=0pt
      #3%
    }%
  }%

  \par
  \nopagebreak[4]%
  \vspace{\wherelistEquationToWhereCorrection}%
  \nopagebreak[4]%
  \noindent where%
  \par\nobreak
  \vspace{\wherelistWhereToItemsCorrection}%
  \nopagebreak[4]%

  % Second pass: typeset the real list using the measured label width.
  \begin{list}{}{%
    \setlength{\labelwidth}{\wherelistLabelWidth}%
    \setlength{\labelsep}{#2}%
    \setlength{\leftmargin}{\dimexpr #1 + \wherelistLabelWidth + #2\relax}%
    \setlength{\itemindent}{0pt}%
    \setlength{\listparindent}{0pt}%
    \setlength{\parsep}{0pt}%
    \setlength{\itemsep}{0pt}%
    \setlength{\topsep}{0pt}%
    \setlength{\partopsep}{0pt}%
    \setlength{\rightmargin}{0pt}%
    \renewcommand{\makelabel}[1]{##1\hfil}%
  }%
  #3%
  \end{list}%

  \endgroup
  \par
  \csname @endpetrue\endcsname
  \ignorespacesafterend
}{}